\paragraph{单原子 Dirichlet 壳、residue 有限部与 cyclotomic 几何}
在单原子因子的周期--primitive 精确剥离、periodic residue 的 Fourier 反演以及 full/core residue 差的 rank-one 局域化已经建立之后，还可把同一原子在 Dirichlet 层继续压缩为一组完全闭式的后果。它们分别把 twisted 通道上的 ghost 贡献识别为精确 polylog 壳，把 residue 类上的 Dirichlet 级数写成 Hurwitz--cyclotomic 壳并统一主极点残数，把 residue-wise 有限部常数的原子位移压缩为显式 digamma 项，并把 centered residue 扰动提升为零均值超平面上的等角紧框架及其精确投影能量律。

\begin{theorem}[single atom 的 twisted Dirichlet 壳是一个精确 polylog 壳]\label{thm:xi-time-part77-atomic-twisted-dirichlet-polylog-shell}
沿用定理 \ref{thm:xi-atomic-witt-cycle-primitive-exact-peeling} 的记号。对固定模数 \(m\ge 2\) 与扭曲通道 \(j\in \ZZ/m\ZZ\)，定义原子 periodic Dirichlet 级数
\[
\mathcal D^{\mathrm{cycle}}_{m,j}(s)
:=
\sum_{n\ge 1}a_n^{\mathrm{cycle}}(\omega_m^j)\,n^{-s}
\qquad (\Re s>1).
\]
则有完全闭式
\[
\mathcal D^{\mathrm{cycle}}_{m,j}(s)
:=
2^{1-s}\operatorname{Li}_s(\omega_m^j),
\]
其中 \(j=0\) 时按 \(\operatorname{Li}_s(1)=\zeta(s)\) 解释。特别地：
\begin{enumerate}
  \item 当 \(j=0\) 时，
  \[
  \mathcal D^{\mathrm{cycle}}_{m,0}(s)=2^{1-s}\zeta(s),
  \]
  因而在 \(s=1\) 处具有简单极点；
  \item 当 \(1\le j\le m-1\) 时，\(\mathcal D^{\mathrm{cycle}}_{m,j}(s)\) 解析延拓为整个函数。
\end{enumerate}
\end{theorem}
\begin{proof}
由定理 \ref{thm:xi-atomic-witt-cycle-primitive-exact-peeling}，
\[
a_n^{\mathrm{cycle}}(u)=
\begin{cases}
2u^{n/2},& 2\mid n,\\
0,& 2\nmid n.
\end{cases}
\]
代入 \(u=\omega_m^j\) 得
\[
a_{2\ell}^{\mathrm{cycle}}(\omega_m^j)=2\omega_m^{j\ell},
\qquad
a_{2\ell+1}^{\mathrm{cycle}}(\omega_m^j)=0.
\]
因此
\[
\mathcal D^{\mathrm{cycle}}_{m,j}(s)
:=
\sum_{\ell\ge 1}2\,\omega_m^{j\ell}(2\ell)^{-s}
:=
2^{1-s}\sum_{\ell\ge 1}\frac{\omega_m^{j\ell}}{\ell^s}
:=
2^{1-s}\operatorname{Li}_s(\omega_m^j),
\]
得到所示闭式。

当 \(j=0\) 时，上式退化为
\[
\mathcal D^{\mathrm{cycle}}_{m,0}(s)=2^{1-s}\zeta(s),
\]
故在 \(s=1\) 处有简单极点。

现设 \(j\neq 0\)。标准 root-of-unity 分解给出
\[
\operatorname{Li}_s(\omega_m^j)
:=
m^{-s}\sum_{r=1}^{m}\omega_m^{jr}\,\zeta\!\left(s,\frac{r}{m}\right).
\]
对每个 \(r\)，Hurwitz zeta \(\zeta(s,r/m)\) 的唯一极点都位于 \(s=1\)，且留数同为 \(1\)。但
\[
\sum_{r=1}^{m}\omega_m^{jr}=0
\qquad (j\neq 0),
\]
故 \(s=1\) 处的极点严格相消，从而 \(\operatorname{Li}_s(\omega_m^j)\) 解析延拓为整个函数。乘上整因子 \(2^{1-s}\) 后，\(\mathcal D^{\mathrm{cycle}}_{m,j}(s)\) 亦为整个函数。证毕。
\end{proof}

\begin{theorem}[single atom 的 periodic residue Dirichlet 级数是 Hurwitz--cyclotomic 壳，且全部 residue 主极点残数完全一致]\label{thm:xi-time-part77-atomic-residue-dirichlet-hurwitz-shell}
沿用定理 \ref{thm:xi-atomic-witt-energy-residue-l2-law} 与结论 \ref{con:xi-time-part57ba-atomic-residue-delta-flat-extremizer} 的记号。对固定模数 \(m\ge 2\) 与 residue 类 \(r\in \ZZ/m\ZZ\)，定义
\[
\mathcal A^{\mathrm{cycle}}_{m,r}(s)
:=
\sum_{n\ge 1}A^{\mathrm{cycle}}_{m,r}(n)\,n^{-s}
\qquad (\Re s>1).
\]
当 \(r\neq 0\) 时，下文把同一符号 \(r\) 也用于其唯一代表元 \(r\in\{1,\dots,m-1\}\)。
则成立：
\begin{enumerate}
  \item 对 \(r=0\)，有
  \[
  \mathcal A^{\mathrm{cycle}}_{m,0}(s)
  =
  2^{1-s}m^{-s}\zeta(s).
  \]
  \item 对 \(1\le r\le m-1\)，有
  \[
  \mathcal A^{\mathrm{cycle}}_{m,r}(s)
  =
  2^{1-s}m^{-s}\zeta\!\left(s,\frac{r}{m}\right).
  \]
\end{enumerate}
特别地，对每个 residue 类 \(r\in \ZZ/m\ZZ\)，\(\mathcal A^{\mathrm{cycle}}_{m,r}(s)\) 在 \(s=1\) 处都有简单极点，并且
\[
\operatorname*{Res}_{s=1}\mathcal A^{\mathrm{cycle}}_{m,r}(s)=\frac{1}{m}.
\]
\end{theorem}
\begin{proof}
由结论 \ref{con:xi-time-part57ba-atomic-residue-delta-flat-extremizer}，
\[
A_{m,r}^{\mathrm{cycle}}(n)
:=
2\,\mathbf 1_{\{2\mid n\}}\mathbf 1_{\{r\equiv n/2\;(\mathrm{mod}\,m)\}}.
\]
若 \(r=0\)，则非零项恰对应于 \(n=2mq\)（\(q\ge 1\)），故
\[
\mathcal A^{\mathrm{cycle}}_{m,0}(s)
:=
2\sum_{q\ge 1}(2mq)^{-s}
:=
2^{1-s}m^{-s}\zeta(s).
\]

若 \(1\le r\le m-1\)，则非零项恰对应于
\[
n=2(r+mq)\qquad (q\ge 0),
\]
故
\[
\mathcal A^{\mathrm{cycle}}_{m,r}(s)
:=
2\sum_{q\ge 0}(2(r+mq))^{-s}
:=
2^{1-s}m^{-s}\sum_{q\ge 0}\left(q+\frac{r}{m}\right)^{-s}
:=
2^{1-s}m^{-s}\zeta\!\left(s,\frac{r}{m}\right).
\]
这就得到前两条公式。

已知 \(\zeta(s)\) 与全部 Hurwitz zeta \(\zeta(s,a)\)（\(0<a\le 1\)）在 \(s=1\) 的留数都等于 \(1\)。于是
\[
\operatorname*{Res}_{s=1}\mathcal A^{\mathrm{cycle}}_{m,r}(s)
:=
\left.2^{1-s}m^{-s}\right|_{s=1}
:=
\frac{1}{m}.
\]
证毕。
\end{proof}

\begin{corollary}[single atom 对 residue-wise Mertens 常数只产生显式 digamma 位移]\label{cor:xi-time-part77-atomic-residue-finitepart-digamma-shift}
对固定模数 \(m\ge 2\) 与 residue 类 \(r\in \ZZ/m\ZZ\)，把 single atom 壳在 \(s=1\) 的 Laurent 展开记为
\[
\mathcal A^{\mathrm{cycle}}_{m,r}(s)
:=
\frac{1/m}{s-1}
+\mathcal M^{\mathrm{cycle}}_{m,r}
+O(s-1).
\]
则
\[
\mathcal M^{\mathrm{cycle}}_{m,0}
:=
\frac{\gamma-\log(2m)}{m},
\]
\[
\mathcal M^{\mathrm{cycle}}_{m,r}
:=
\frac{-\psi(r/m)-\log(2m)}{m}
\qquad (1\le r\le m-1),
\]
其中 \(\gamma\) 为 Euler 常数，\(\psi=\Gamma'/\Gamma\) 为 digamma 函数。

进一步，定义 full/core residue Dirichlet 级数
\[
\mathcal A^{\mathrm{full}}_{m,r}(s)
:=
\sum_{n\ge 1}A^{\mathrm{full}}_{m,r}(n)\,n^{-s},
\qquad
\mathcal A^{\mathrm{core}}_{m,r}(s)
:=
\sum_{n\ge 1}A^{\mathrm{core}}_{m,r}(n)\,n^{-s}.
\]
若把相应 Laurent 常数项记为
\[
\mathcal A^{\star}_{m,r}(s)
:=
\frac{\kappa^{\star}_{m,r}}{s-1}
+\log\mathfrak M^{\star}_{m,r}
+O(s-1),
\qquad
\star\in\{\mathrm{full},\mathrm{core}\},
\]
并在 \(r\neq 0\) 时仍以 \(r\in\{1,\dots,m-1\}\) 表示其标准代表元，则有
\[
\log\mathfrak M^{\mathrm{full}}_{m,0}
-\log\mathfrak M^{\mathrm{core}}_{m,0}
:=
\mathcal M^{\mathrm{cycle}}_{m,0}
=
\frac{\gamma-\log(2m)}{m},
\]
\[
\log\mathfrak M^{\mathrm{full}}_{m,r}
-\log\mathfrak M^{\mathrm{core}}_{m,r}
:=
\mathcal M^{\mathrm{cycle}}_{m,r}
=
\frac{-\psi(r/m)-\log(2m)}{m}
\qquad (1\le r\le m-1).
\]
因此，single atom 在 residue-wise finite part 层并不与 core 混合，而只留下由 cyclotomic/Hurwitz 数据完全决定的一族 digamma 位移。
\end{corollary}
\begin{proof}
由定理 \ref{thm:xi-time-part77-atomic-residue-dirichlet-hurwitz-shell} 的闭式表达式与标准展开
\[
\zeta(s)=\frac{1}{s-1}+\gamma+O(s-1),
\qquad
\zeta(s,a)=\frac{1}{s-1}-\psi(a)+O(s-1),
\]
以及
\[
2^{1-s}m^{-s}
=
\frac{1}{m}\Bigl(1-(s-1)\log(2m)+O((s-1)^2)\Bigr),
\]
逐项相乘，即得 \(\mathcal M^{\mathrm{cycle}}_{m,0}\) 与 \(\mathcal M^{\mathrm{cycle}}_{m,r}\) 的两条显式公式。

另一方面，由定义，
\[
A^{\mathrm{full}}_{m,r}(n)=A^{\mathrm{core}}_{m,r}(n)+A^{\mathrm{cycle}}_{m,r}(n),
\]
故在 \(\Re s>1\) 上
\[
\mathcal A^{\mathrm{full}}_{m,r}(s)
:=
\mathcal A^{\mathrm{core}}_{m,r}(s)+\mathcal A^{\mathrm{cycle}}_{m,r}(s).
\]
于是 \(\log\mathfrak M^{\mathrm{full}}_{m,r}-\log\mathfrak M^{\mathrm{core}}_{m,r}\) 恰等于 \(\mathcal A^{\mathrm{cycle}}_{m,r}(s)\) 在 \(s=1\) 处的 Laurent 常数项，即 \(\mathcal M^{\mathrm{cycle}}_{m,r}\)。证毕。
\end{proof}

\begin{theorem}[single atom 的 centered residue 扰动在零均值超平面上构成等角紧框架]\label{thm:xi-time-part77-atomic-centered-residue-tight-frame}\leanverified{paper\_xi\_time\_part77\_atomic\_centered\_residue\_tight\_frame}
沿用定理 \ref{thm:xi-time-part65c-atomic-centered-residue-simplex} 的记号。对固定模数 \(m\ge 2\)，令
\[
\mathcal H_m
:=
\left\{x\in\RR^{\ZZ/m\ZZ}:\ \sum_{r\in\ZZ/m\ZZ}x_r=0\right\}.
\]
则族 \(\{\Delta^{(\ell)}\}_{\ell\in\ZZ/m\ZZ}\) 满足：
\begin{enumerate}
  \item 对任意 \(\ell,\ell'\in\ZZ/m\ZZ\)，有 \(\Delta^{(\ell)}\in\mathcal H_m\) 且
  \[
  \bigl\langle \Delta^{(\ell)},\Delta^{(\ell')}\bigr\rangle
  =
  4\left(\delta_{\ell,\ell'}-\frac{1}{m}\right).
  \]
  \item 因而 \(\{\Delta^{(\ell)}\}_{\ell\in\ZZ/m\ZZ}\) 在 \(\mathcal H_m\) 中构成一个等角紧框架；更精确地，
  \[
  \frac{1}{m}\sum_{\ell\in\ZZ/m\ZZ}
  \Delta^{(\ell)}\bigl(\Delta^{(\ell)}\bigr)^{\mathsf T}
  =
  \frac{4}{m}\left(I-\frac{1}{m}J\right),
  \]
  其中 \(J\) 为全 \(1\) 矩阵。
  \item 对任意子空间 \(V\subset\mathcal H_m\) 及其正交投影 \(Q_V\)，若 \(\dim V=d\)，则
  \[
  \frac{1}{m}\sum_{\ell\in\ZZ/m\ZZ}
  \bigl\|Q_V\Delta^{(\ell)}\bigr\|_2^2
  =
  \frac{4d}{m}.
  \]
\end{enumerate}
换言之，single atom 的 centered residue shell 在零均值超平面上对任意 \(d\)-维线性探针的平均能量注入都只由维数 \(d\) 决定，而与探针方向无关。
\end{theorem}
\begin{proof}
第 \(1\) 条正是定理 \ref{thm:xi-time-part65c-atomic-centered-residue-simplex} 的内容。又由推论 \ref{cor:xi-time-part65c-atomic-centered-residue-isotropic}，
\[
\frac{1}{m}\sum_{\ell\in\ZZ/m\ZZ}
\Delta^{(\ell)}\bigl(\Delta^{(\ell)}\bigr)^{\mathsf T}
=
\frac{4}{m}\left(I-\frac{1}{m}J\right).
\]
其中 \(I-\frac{1}{m}J\) 正是到零均值超平面 \(\mathcal H_m\) 的正交投影，故第 \(2\) 条给出 \(\{\Delta^{(\ell)}\}\) 在 \(\mathcal H_m\) 上的紧框架性质；而第 \(1\) 条中的两两内积公式又表明该框架是等角的。

最后，对任意 \(V\subset\mathcal H_m\) 的正交投影 \(Q_V\) 取迹可得
\[
\frac{1}{m}\sum_{\ell\in\ZZ/m\ZZ}
\bigl\|Q_V\Delta^{(\ell)}\bigr\|_2^2
=
\operatorname{tr}\!\left(
Q_V\cdot
\frac{1}{m}\sum_{\ell\in\ZZ/m\ZZ}
\Delta^{(\ell)}\bigl(\Delta^{(\ell)}\bigr)^{\mathsf T}
\right)
=
\frac{4}{m}\operatorname{tr}\!\left(Q_V\left(I-\frac{1}{m}J\right)\right).
\]
由于 \(V\subset\mathcal H_m=\mathbf 1^\perp\)，故 \(Q_V\mathbf 1=0\)，从而 \(Q_VJ=0\)。因此
\[
\operatorname{tr}\!\left(Q_V\left(I-\frac{1}{m}J\right)\right)
=
\operatorname{tr}(Q_V)
=
d,
\]
便得到
\[
\frac{1}{m}\sum_{\ell\in\ZZ/m\ZZ}
\bigl\|Q_V\Delta^{(\ell)}\bigr\|_2^2
=
\frac{4d}{m}.
\]
证毕。
\end{proof}

\noindent
这四条结果把单原子因子的 residue / twisted 接口进一步压缩到纯 Dirichlet 层的完全显式外壳：在 twisted 通道上，它恰等于一个根单位 polylog 壳；在 residue 通道上，它恰等于一族 Hurwitz--cyclotomic 壳，并把全部主极点残数统一锁定为 \(1/m\)；相应地，full/core residue-wise finite part 常数之间的差也不再携带任何隐藏动力学信息，而只是一组显式的 digamma 位移；最后，centered residue 扰动不仅给出正则单纯形几何，而且在零均值超平面上形成等角紧框架，使任意低维线性探针所接收到的平均能量严格服从维数比例律。

\paragraph{相位采样 \texorpdfstring{$\zeta$}{zeta} 的 Hurwitz 闭式、群刚性重建与异常分裂 torsor 势函数}
沿用定义 \ref{def:cdim-zeta-function}、命题 \ref{prop:cdim-residual-quotient-probe}、定理 \ref{thm:cdim-anomaly-quadratic-law}、命题 \ref{prop:cdim-discrete-anomaly-ledger} 与命题 \ref{prop:cdim-anomaly-splitting-torsor-obstruction} 的记号。对有限生成离散交换群
\[
G\cong \ZZ^r\oplus T
\]
记
\[
a_n(G):=\card{\operatorname{Hom}(G,\ZZ/n\ZZ)},
\qquad
\zeta_G^{\mathrm{ph}}(s):=\sum_{n\ge 1}\frac{a_n(G)}{n^s},
\]
并把 \(\zeta_G^{\mathrm{ph}}\) 视为前文 \(\zeta_G\) 的相位采样记号。由有限采样计数、异常群分解与分裂截面的 torsor 结构，可进一步压缩出下列五条结论。

\begin{theorem}[相位采样 \texorpdfstring{$\zeta$}{zeta} 的 Hurwitz 分解与唯一极点]\label{thm:xi-time-part78-phase-zeta-hurwitz-pole}
\leanverified{paper\_xi\_time\_part78\_phase\_zeta\_hurwitz\_pole}
设 \(G\cong \ZZ^r\oplus T\) 为有限生成离散交换群，其中 \(T\) 有限。记
\[
e:=\exp(T),
\qquad
c_T(n):=\card{T/nT}.
\]
则成立：
\begin{enumerate}
  \item \(c_T(n)\) 仅依赖于 \(\gcd(n,e)\)，因而是以 \(e\) 为周期的算术函数；
  \item 在 \(\Re(s)>r+1\) 上有 Hurwitz 分解
  \[
  \zeta_G^{\mathrm{ph}}(s)
  =e^{\,r-s}\sum_{a=1}^{e}c_T(a)\,\zeta\!\left(s-r,\frac{a}{e}\right);
  \]
  \item \(\zeta_G^{\mathrm{ph}}(s)\) 亚纯延拓到整个复平面，并在 \(s=r+1\) 处有唯一单极点，其留数为
  \[
  \operatorname*{Res}_{s=r+1}\zeta_G^{\mathrm{ph}}(s)
  =\frac{1}{e}\sum_{a=1}^{e}c_T(a).
  \]
\end{enumerate}
\end{theorem}
\begin{proof}
由命题 \ref{prop:cdim-residual-quotient-probe}，
\[
a_n(G)=n^r\card{T/nT}=n^r c_T(n).
\]
将 \(T\) 分解为 \(p\)-初等分量 \(T=\bigoplus_p T_{(p)}\)。对任意素数 \(p\) 与任意整数 \(n\)，乘以 \(n\) 在 \(T_{(p)}\) 上可写成
\[
n=p^{v_p(n)}u,\qquad p\nmid u,
\]
而乘以 \(u\) 为自同构，故
\[
\card{T_{(p)}/nT_{(p)}}=\card{T_{(p)}/p^{v_p(n)}T_{(p)}}.
\]
若 \(p^{\lambda_p}\Vert e\)，则 \(T_{(p)}\) 的指数为 \(p^{\lambda_p}\)，从而上式只依赖于 \(\min(v_p(n),\lambda_p)\)，亦即只依赖于 \(\gcd(n,p^{\lambda_p})\)。对所有 \(p\mid e\) 取乘积便得 \(c_T(n)\) 只依赖于 \(\gcd(n,e)\)，故以 \(e\) 为周期。

于是当 \(\Re(s)>r+1\) 时，
\[
\zeta_G^{\mathrm{ph}}(s)
=\sum_{n\ge 1}c_T(n)\,n^{r-s}
=\sum_{a=1}^{e}c_T(a)\sum_{m\ge 0}(me+a)^{r-s}.
\]
将 \(me+a=e\bigl(m+\frac{a}{e}\bigr)\) 提出，得到
\[
\zeta_G^{\mathrm{ph}}(s)
=e^{\,r-s}\sum_{a=1}^{e}c_T(a)\sum_{m\ge 0}\left(m+\frac{a}{e}\right)^{-(s-r)}
=e^{\,r-s}\sum_{a=1}^{e}c_T(a)\,\zeta\!\left(s-r,\frac{a}{e}\right).
\]
这便给出所述 Hurwitz 分解。

每个 Hurwitz zeta \(\zeta(s-r,a/e)\) 都只在 \(s-r=1\) 处有单极点，留数为 \(1\)。因此 \(\zeta_G^{\mathrm{ph}}\) 亚纯延拓到整个复平面，并且只能在 \(s=r+1\) 处出现极点；其留数为
\[
e^{\,r-(r+1)}\sum_{a=1}^{e}c_T(a)=\frac1e\sum_{a=1}^{e}c_T(a).
\]
由于 \(c_T(a)\ge 1\) 对所有 \(a\) 都成立，该留数严格为正，故该极点确为唯一单极点。证毕。
\end{proof}

\begin{theorem}[相位采样计数部分和的精确主项与 \texorpdfstring{$O(N^r)$}{O(N^r)} 余项]\label{thm:xi-time-part78-phase-zeta-partial-sum-sharp}
沿用定理 \ref{thm:xi-time-part78-phase-zeta-hurwitz-pole} 的记号，并记
\[
A_G(N):=\sum_{n\le N}a_n(G).
\]
则当 \(N\to\infty\) 时，
\[
A_G(N)=\frac{C_G}{r+1}N^{r+1}+O(N^r),
\qquad
C_G:=\frac1e\sum_{a=1}^{e}c_T(a).
\]
特别地，
\[
C_G=\operatorname*{Res}_{s=r+1}\zeta_G^{\mathrm{ph}}(s).
\]
\end{theorem}
\begin{proof}
由定理 \ref{thm:xi-time-part78-phase-zeta-hurwitz-pole}，函数 \(c_T(n)\) 以 \(e\) 为周期，故
\[
A_G(N)
=\sum_{a=1}^{e}c_T(a)\sum_{\substack{n\le N\\ n\equiv a\!\!\!\pmod e}}n^r.
\]
对每个 \(a\in\{1,\dots,e\}\) 记
\[
M_a:=\left\lfloor\frac{N-a}{e}\right\rfloor.
\]
则内层和可写为
\[
\sum_{\substack{n\le N\\ n\equiv a\!\!\!\pmod e}}n^r
=\sum_{m=0}^{M_a}(em+a)^r.
\]
由于 \((em+a)^r\) 是 \(m\) 的 \(r\) 次多项式，Faulhaber 求和公式给出
\[
\sum_{m=0}^{M_a}(em+a)^r
=\frac{e^r}{r+1}M_a^{r+1}+O(M_a^r),
\]
其中隐含常数对 \(a=1,\dots,e\) 一致，因为 \(e\) 固定且 \(a\) 只在有限集合中变动。又
\[
M_a=\frac{N}{e}+O(1)
\]
同样对 \(a\) 一致，于是
\[
\sum_{m=0}^{M_a}(em+a)^r
=\frac{1}{e(r+1)}N^{r+1}+O(N^r).
\]
代回并对 \(a\) 求和，得到
\[
A_G(N)
=\sum_{a=1}^{e}c_T(a)\left(\frac{1}{e(r+1)}N^{r+1}+O(N^r)\right)
=\frac{C_G}{r+1}N^{r+1}+O(N^r).
\]
最后，\(C_G=\operatorname*{Res}_{s=r+1}\zeta_G^{\mathrm{ph}}(s)\) 正是定理 \ref{thm:xi-time-part78-phase-zeta-hurwitz-pole} 的留数公式。证毕。
\end{proof}

\begin{theorem}[相位采样 \texorpdfstring{$\zeta$}{zeta} 的完全分类刚性]\label{thm:xi-time-part78-phase-zeta-complete-classification}
设 \(G,H\) 为有限生成离散交换群。若
\[
\zeta_G^{\mathrm{ph}}(s)\equiv \zeta_H^{\mathrm{ph}}(s)
\]
作为亚纯函数恒等成立，则
\[
G\cong H.
\]
等价地，映射
\[
G\longmapsto \zeta_G^{\mathrm{ph}}(s)
\]
是有限生成离散交换群同构类的完备不变量。
\end{theorem}
\begin{proof}
在某个足够右的半平面内，两端都由绝对收敛的 Dirichlet 级数
\[
\zeta_G^{\mathrm{ph}}(s)=\sum_{n\ge 1}\frac{a_n(G)}{n^s},
\qquad
\zeta_H^{\mathrm{ph}}(s)=\sum_{n\ge 1}\frac{a_n(H)}{n^s}
\]
给出。亚纯恒等因而在该半平面内也成立。由 Dirichlet 级数唯一性定理，
\[
a_n(G)=a_n(H)\qquad(\forall n\ge 1).
\]
再由推论 \ref{cor:xi-cdim-zeta-fingerprint-rigidity}（等价地，定理 \ref{thm:cdim-phase-spectrum-reconstruction}），可知这迫使 \(G\cong H\)。证毕。
\end{proof}

\begin{theorem}[异常分裂 torsor 的精确势函数]\label{thm:xi-time-part78-anomaly-splitting-torsor-cardinality}
设
\[
G\cong \ZZ^r\oplus F
\]
为有限生成离散交换群，其中 \(F\) 有限。则分裂截面集合 \(\mathrm{Spl}(G)\) 为有限集，并满足
\[
\card{\mathrm{Spl}(G)}
=\bigl(\card{F}^{\,r}\,\card{\wedge^2 F}\bigr)^{\binom r2}.
\]
更具体地，若对每个素数 \(p\) 写
\[
F_{(p)}\cong \bigoplus_{j=1}^{t_p}\ZZ/p^{\lambda_{p,j}}\ZZ,
\qquad
\lambda_{p,1}\ge \cdots\ge \lambda_{p,t_p}\ge 1,
\]
则
\[
\card{\wedge^2 F}
=\prod_p p^{\sum_{j=2}^{t_p}(j-1)\lambda_{p,j}},
\]
从而
\[
\card{\mathrm{Spl}(G)}
=\prod_p
p^{\binom r2\left(
r\sum_{j=1}^{t_p}\lambda_{p,j}
+\sum_{j=2}^{t_p}(j-1)\lambda_{p,j}
\right)}.
\]
特别地，当 \(r\le 1\) 或 \(F=0\) 时，\(\mathrm{Spl}(G)\) 仅含一个元素。
\end{theorem}
\begin{proof}
由命题 \ref{prop:cdim-anomaly-splitting-torsor-obstruction}，
\(\mathrm{Spl}(G)\) 是
\[
\operatorname{Hom}\bigl(\mathcal D(G),\mathcal A(G)^0\bigr)
\]
的主齐次空间，因而
\[
\card{\mathrm{Spl}(G)}
=\card{\operatorname{Hom}\bigl(\mathcal D(G),\mathcal A(G)^0\bigr)}.
\]
再由定理 \ref{thm:cdim-anomaly-quadratic-law}，
\[
\mathcal A(G)^0\cong \TT^{\binom r2},
\]
故
\[
\operatorname{Hom}\bigl(\mathcal D(G),\mathcal A(G)^0\bigr)
\cong
\operatorname{Hom}\bigl(\mathcal D(G),\TT\bigr)^{\binom r2}.
\]
由于 \(\mathcal D(G)\) 为有限交换群，Pontryagin 对偶给出
\[
\card{\operatorname{Hom}\bigl(\mathcal D(G),\TT\bigr)}=\card{\mathcal D(G)},
\]
从而
\[
\card{\mathrm{Spl}(G)}=\card{\mathcal D(G)}^{\binom r2}.
\]

另一方面，命题 \ref{prop:cdim-discrete-anomaly-ledger} 给出
\[
\mathcal D(G)\cong
\operatorname{Hom}(\ZZ^r\otimes F,\TT)\times \operatorname{Hom}(\wedge^2F,\TT).
\]
因此
\[
\card{\mathcal D(G)}
=\card{\ZZ^r\otimes F}\,\card{\wedge^2F}
=\card{F}^{\,r}\,\card{\wedge^2F},
\]
因为 \(\ZZ^r\otimes F\cong F^r\)。于是得到
\[
\card{\mathrm{Spl}(G)}
=\bigl(\card{F}^{\,r}\,\card{\wedge^2F}\bigr)^{\binom r2}.
\]

为求 \(\card{\wedge^2F}\)，先按主分解写 \(F=\bigoplus_p F_{(p)}\)。由不同素数方向的张量积消失，
\[
\wedge^2F\cong \bigoplus_p \wedge^2F_{(p)}.
\]
再对
\[
F_{(p)}\cong \bigoplus_{j=1}^{t_p}\ZZ/p^{\lambda_{p,j}}\ZZ
\]
应用有限交换 \(p\)-群的外平方公式，得到
\[
\wedge^2F_{(p)}
\cong
\bigoplus_{1\le i<j\le t_p}\ZZ/p^{\min(\lambda_{p,i},\lambda_{p,j})}\ZZ.
\]
由于 \(\lambda_{p,1}\ge\cdots\ge\lambda_{p,t_p}\)，当 \(i<j\) 时有
\(\min(\lambda_{p,i},\lambda_{p,j})=\lambda_{p,j}\)，故
\[
\card{\wedge^2F_{(p)}}
=p^{\sum_{j=2}^{t_p}(j-1)\lambda_{p,j}}.
\]
对所有 \(p\) 取乘积便得
\[
\card{\wedge^2F}
=\prod_p p^{\sum_{j=2}^{t_p}(j-1)\lambda_{p,j}}.
\]
再与
\[
\card{F}=\prod_p p^{\sum_{j=1}^{t_p}\lambda_{p,j}}
\]
合并，即得最后的素数分解式。

若 \(r\le 1\)，则 \(\binom r2=0\)，故 \(\card{\mathrm{Spl}(G)}=1\)。若 \(F=0\)，则 \(\mathcal D(G)=0\)，同样得到 \(\mathrm{Spl}(G)\) 为单点。证毕。
\end{proof}

\begin{corollary}[异常分裂 torsor 的素数支撑与严格下界]\label{cor:xi-time-part78-anomaly-splitting-prime-support-lower-bound}
沿用定理 \ref{thm:xi-time-part78-anomaly-splitting-torsor-cardinality} 的记号，并设 \(r\ge 2\)。则对任意素数 \(p\) 都有
\[
p\mid \card{\mathrm{Spl}(G)}
\quad\Longleftrightarrow\quad
F_{(p)}\neq 0.
\]
等价地，
\[
\{p:\ p\mid \card{\mathrm{Spl}(G)}\}
=
\{p:\ F_{(p)}\neq 0\}.
\]
若进一步 \(F\neq 0\)，并记
\[
p_{\min}(F):=\min\{p:\ F_{(p)}\neq 0\},
\]
则
\[
\card{\mathrm{Spl}(G)}
\ge p_{\min}(F)^{\,r\binom r2}.
\]
\end{corollary}
\begin{proof}
由定理 \ref{thm:xi-time-part78-anomaly-splitting-torsor-cardinality} 的显式分解式，
\[
v_p\!\bigl(\card{\mathrm{Spl}(G)}\bigr)
=\binom r2\left(
r\sum_{j=1}^{t_p}\lambda_{p,j}
+\sum_{j=2}^{t_p}(j-1)\lambda_{p,j}
\right).
\]
若 \(F_{(p)}=0\)，则所有 \(\lambda_{p,j}\) 均为零，从而
\(
v_p(\card{\mathrm{Spl}(G)})=0
\)。
反之若 \(F_{(p)}\neq 0\)，则
\(
\sum_{j=1}^{t_p}\lambda_{p,j}\ge 1
\)；
又因 \(r\ge 2\)，故
\[
v_p\!\bigl(\card{\mathrm{Spl}(G)}\bigr)
\ge \binom r2\,r>0.
\]
于是
\[
p\mid \card{\mathrm{Spl}(G)}
\quad\Longleftrightarrow\quad
F_{(p)}\neq 0.
\]

若 \(F\neq 0\)，取其最小支撑素数 \(p_{\min}(F)\)。对该素数有
\[
v_{p_{\min}(F)}\!\bigl(\card{\mathrm{Spl}(G)}\bigr)\ge \binom r2\,r,
\]
其余素数因子贡献非负，故
\[
\card{\mathrm{Spl}(G)}
\ge p_{\min}(F)^{\,r\binom r2}.
\]
证毕。
\end{proof}

\noindent
这五条结果把有限生成离散交换群的 phase-sampling 侧统计与二阶异常分裂的离散扭量压缩到同一条可审计主线上：相位采样 \(\zeta\) 不仅具有有限 Hurwitz 闭式并把唯一主极点精确锁定在 \(s=\cdim(G)+1\)，而且其亚纯函数本身已经足以完整恢复群的同构类型；与此同时，异常分裂截面的全部非唯一性也不再只是存在性障碍，而被一个可显式计算的有限 torsor 势函数完全控制，其素数支撑与大小下界均由 torsion 主分解直接给出。

\paragraph{MOM 乘法分级、交叉异常阈值与 window-\texorpdfstring{$6$}{6} 扇区读出的信息压缩}
在相位采样 \(\zeta\) 的 Hurwitz 闭式、群刚性重建以及异常分裂 torsor 的势函数已经闭合之后，前文关于四生成元投影词重写、交叉异常坐标化、对称幂 moment-kernel 编译与 window-\(6\) 的取向扇区结构还可继续压缩为同一条更高层的审计链：MOM 前缀指数给出乘法分级，交叉异常族给出轴可分解性的最小完备审计坐标，而对称幂编译与四扇区读出则分别把“显影成本”与“信息损失”严格量化。

\begin{theorem}[四生成元投影词演算的乘法分级分裂]\label{thm:xi-time-part79-pom-multiplicative-graded-splitting}
在定义 \ref{def:pom-four-gen-rewrite-rules}、定义 \ref{def:pom-algorithmic-equivalence-nf-anom} 与命题 \ref{prop:pom-confluence-mod-anom-audited} 的有界审计切片内，令
\[
\mathcal W_4
\]
为由
\(
\mathrm{PROJ}_u,\ E_m,\ \mathrm{LIFT}_G,\ \mathrm{MOM}_q\ (q\ge 2)
\)
生成的自由单半群。记
\[
\Pi_{\mathrm{MOM}}:\mathcal W_4\to (\NN_{\ge 1},\times)
\]
为命题 \ref{prop:pom-mom-prefix-product-invariant} 的 MOM 前缀指数。称 \(w_1,w_2\in\mathcal W_4\) 结构等价，记为
\(
w_1\equiv_{\mathrm{str}}w_2
\)，若二者都可仅用 \((\mathrm{RMOM})\) 与 \((\mathrm{RMOM\text{-}BUBBLE})\) 规约为
\[
w_i\ \Rightarrow\ \mathrm{MOM}_{Q_i}\circ v_i,
\qquad
v_i\ \text{不含}\ \mathrm{MOM},
\]
并满足
\[
Q_1=Q_2
\qquad\text{且}\qquad
v_1\sim v_2,
\]
其中 \(\sim\) 为定义 \ref{def:pom-algorithmic-equivalence-nf-anom} 的三生成元算法等价。
则存在规范分裂
\[
\mathcal W_4/\equiv_{\mathrm{str}}
\ \cong\
\bigsqcup_{Q\ge 1}\mathfrak E_Q,
\qquad
\mathfrak E_Q:=
\{[w]_{\equiv_{\mathrm{str}}}:\ \Pi_{\mathrm{MOM}}(w)=Q\}.
\]
并且每个结构类都存在一个前缀正规代表
\[
\mathrm{MOM}_{Q}\circ N,
\]
其中 \(N\) 为命题 \ref{prop:pom-three-gen-interface-normal-form} 的三生成元接口正规形。更精确地，在固定 \(Q\) 的扇区 \(\mathfrak E_Q\) 内，结构类恰由不含 MOM 的残余片段的
\[
\bigl(\mathrm{NF},[\mathrm{Anom}]\bigr)
\]
唯一决定。
\end{theorem}
\begin{proof}
命题 \ref{prop:pom-mom-prefix-product-invariant} 断言 \(\Pi_{\mathrm{MOM}}\) 在允许重写下严格不变，因此不同 \(Q\) 的词不可能在结构等价中相遇。
另一方面，命题 \ref{prop:pom-mom-interface-normal-form} 给出一条规范的归一化策略：先用 \((\mathrm{RMOM})\) 合并全部 moment 门，再用 \((\mathrm{RMOM\text{-}BUBBLE})\) 将其外提为唯一前缀 \(\mathrm{MOM}_Q\)，从而把任意 \(w\) 化为 \(\mathrm{MOM}_Q\circ v\) 的形式，其中 \(v\) 落回三生成元片段。
对该不含 MOM 的残余片段，命题 \ref{prop:pom-confluence-mod-anom-audited} 与定义 \ref{def:pom-algorithmic-equivalence-nf-anom} 表明，其路径不变量正是接口正规形与异常上同调类；因此在固定 \(Q\) 内，结构类恰由 \((\mathrm{NF}(v),[\mathrm{Anom}(v)])\) 分类。综上，整个商集按 MOM 前缀指数分裂为互不连通的乘法分级扇区。证毕。
\end{proof}

\begin{theorem}[高阶交叉异常审计的锐完备阈值]\label{thm:xi-time-part79-cross-anom-sharp-completeness-threshold}
\leanverified{paper\_xi\_time\_part79\_cross\_anom\_sharp\_completeness\_threshold}
设 \(k\ge 3\)，并令 \(G_1,\dots,G_k\) 为有限阿贝尔群。固定 \(\theta\) 后，把
\[
f(\chi):=\log\mathfrak{M}_{\chi}(\theta),
\qquad
\chi\in \widehat G_1\times\cdots\times\widehat G_k,
\]
视为常数层载荷，并记其 Möbius 交叉分量为 \(\Delta_S(\chi_S;\theta)\)。
则全体高阶交叉异常
\[
\{\Delta_S:\ |S|\ge 2\}
\]
构成轴可分解性的最小完备审计家族，具体地：
\begin{enumerate}
  \item \(f\) 轴可分解当且仅当对一切 \(|S|\ge 2\) 都有
  \[
  \Delta_S(\chi_S;\theta)\equiv 0.
  \]
  \item 对任意固定的 \(r\in\{3,\dots,k\}\)，只要前 \(r\) 条轴各含一个非平凡角色，则任何仅检查全部
  \[
  \{\Delta_S:\ 2\le |S|<r\}
  \]
  的审计协议都不是完备的。更强地，存在一个纯 \(r\)-体载荷 \(f_r\)，使得
  \[
  \Delta_S^{(r)}\equiv 0\qquad(2\le |S|<r),
  \qquad
  \Delta_{[r]}^{(r)}\not\equiv 0.
  \]
\end{enumerate}
因此，一旦省略任意阶数 \(r\ge 3\) 的检测，就会留下纯 \(r\)-体异常的结构性盲区。
\end{theorem}
\begin{proof}
第 (1) 条正是定理 \ref{thm:pom-axis-decomposability-iff} 的内容。

对第 (2) 条，取前 \(r\) 个轴上的非平凡角色 \(a_i\in\widehat G_i\setminus\{\chi_{i,0}\}\)，并定义
\[
f_r(\chi):=
\begin{cases}
1,& \chi=(a_1,\dots,a_r,\chi_{r+1,0},\dots,\chi_{k,0}),\\
0,& \text{其他情形}.
\end{cases}
\]
若 \(S\subseteq[k]\) 满足 \(2\le |S|<r\)，则在 \(\Delta_S\) 的 Möbius 差分中，每一项都至少有一个前 \(r\) 坐标被置为平凡角色，故所有取值均为 \(0\)，从而
\(
\Delta_S^{(r)}\equiv 0
\)。
但当 \(S=[r]\) 且 \(\chi_{[r]}=(a_1,\dots,a_r)\) 时，差分和中只有满集项保留，故
\[
\Delta_{[r]}^{(r)}(a_1,\dots,a_r;\theta)=1\neq 0.
\]
于是 \(f_r\) 不是轴可分解的，却逃过了所有低于 \(r\) 阶的交叉异常审计。证毕。
\end{proof}

\begin{theorem}[异常显影阈值的多项式编译成本]\label{thm:xi-time-part79-anom-threshold-polynomial-compilation-cost}
设 \(w\) 为任意投影词，并记
\[
a:=[\mathrm{Anom}(w)]
\]
为其异常上同调类。若 \(a\neq 0\)，并且底层在线核的内部状态集为有限集 \(\mathcal Q\)，则对任意阈值 \(\tau>0\) 都存在整数 \(q\ge 2\)，使得：
\begin{enumerate}
  \item 某个异常坐标 \(\chi\) 满足
  \[
  \bigl|[\mathrm{Anom}(w^{\boxtimes q})]_{\chi}\bigr|\ge \tau;
  \]
  \item \(q\)-重 moment-kernel 可在命题 \ref{prop:pom-moment-symmetric-power} 的对称幂商核上实现，其状态数至多为
  \[
  \binom{q+|\mathcal Q|-1}{|\mathcal Q|-1};
  \]
  \item 因而实现阈值 \(\tau\) 的最小状态成本满足上界
  \[
  \mathrm{Cost}_{\min}(\tau)=O\!\bigl(\tau^{|\mathcal Q|-1}\bigr).
  \]
\end{enumerate}
与此相对，若采用未压缩的 \(q\)-份同步直积实现，则状态数为 \(|\mathcal Q|^q\)，即指数规模。
\end{theorem}
\begin{proof}
由命题 \ref{prop:pom-anom-gap-amplification}，有
\[
[\mathrm{Anom}(w^{\boxtimes q})]=q\cdot a.
\]
取一坐标 \(\chi\) 使 \(a_\chi\neq 0\)，并令
\[
q:=\max\!\left\{2,\left\lceil \frac{\tau}{|a_\chi|}\right\rceil\right\}.
\]
则
\[
\bigl|[\mathrm{Anom}(w^{\boxtimes q})]_{\chi}\bigr|
=
q|a_\chi|
\ge \tau.
\]
另一方面，命题 \ref{prop:pom-moment-symmetric-power} 给出 \(q\)-重碰撞核的对称幂商态空间
\[
\mathcal N_q=\left\{(n_s)_{s\in\mathcal Q}\in\ZZ_{\ge 0}^{\mathcal Q}:\ \sum_{s\in\mathcal Q}n_s=q\right\},
\]
其基数正是
\[
|\mathcal N_q|=\binom{q+|\mathcal Q|-1}{|\mathcal Q|-1}.
\]
由于 \(q=O(\tau)\)，故最小状态成本至多按 \(\tau^{|\mathcal Q|-1}\) 多项式增长。若不做对称压缩，则直接同步直积的状态空间为 \(\mathcal Q^q\)，大小为 \(|\mathcal Q|^q\)，故为指数级。证毕。
\end{proof}

\begin{theorem}[window-\texorpdfstring{$6$}{6} 取向扇区读出的 \texorpdfstring{$19$}{19} 比特压缩下界]\label{thm:xi-time-part79-window6-sector-19bit-collapse}
沿用推论 \ref{cor:xi-window6-boundary-parity-unique-minimal-faithful-quotient} 的记号，记
\[
\mathcal A_6:=\bigl(G_6^{\mathrm{bin}}\bigr)^{\mathrm{ab}}\cong (\ZZ/2\ZZ)^{21}
\]
为 window-\(6\) 的完备二元异常签名空间。
另一方面，由命题 \ref{prop:fold-groupoid-z2x2-sector-by-orientation} 与定理 \ref{thm:op-algebra-fold-weak-hopf-z2x2-sector-closure}，window-\(6\) 的取向指纹只给出四个弱 Hopf 扇区
\[
\CC[\mathcal G_6]\cong \bigoplus_{(\alpha,\beta)\in\{\pm 1\}^2} H_6^{\alpha,\beta}.
\]
则任意只通过取向扇区标签
\[
\sigma:\mathcal A_6\to \{\pm 1\}^2
\]
进行的异常分类，都至少丢失 \(19\) 比特信息。更精确地，至少存在一个扇区，其纤维满足
\[
\#\sigma^{-1}(\alpha,\beta)\ge 2^{19}.
\]
\end{theorem}
\begin{proof}
由 \(\mathcal A_6\cong(\ZZ/2\ZZ)^{21}\)，可知
\[
|\mathcal A_6|=2^{21}.
\]
而取向扇区标签集 \(\{\pm1\}^2\) 的基数为 \(4\)。
因此任何 \(\sigma:\mathcal A_6\to\{\pm1\}^2\) 的像至多有 \(4\) 个点。由鸽巢原理，至少有一个标签的原像大小不小于
\[
\frac{2^{21}}{4}=2^{19}.
\]
换言之，取向扇区读出至多保留 \(2\) 个二元比特，而完整异常签名需要 \(21\) 个独立比特；两者之间的不可约信息缺口至少为 \(19\) 比特。证毕。
\end{proof}

\begin{conjecture}[取向扇区读出的线性失配律]\label{conj:xi-time-part79-orientation-sector-linear-mismatch}
设 \(m\) 满足所有 bin-fold 纤维都非平凡，并令 \(\mathcal A_m\) 为保持纤维独立置换语义的最小完备二元异常签名空间。则任意仅通过取向指纹的四扇区标签
\[
\sigma_m:\mathcal A_m\to \{\pm 1\}^2
\]
实现的异常读出，在所有与弱 Hopf 扇区分裂相容的实现中，都应满足不可改进的信息损失下界
\[
\mathrm{loss}(\sigma_m)\ge |X_m|-2.
\]
换言之，四扇区数保持常值，而完备异常维数应随分辨率按 \(|X_m|\) 线性增长；window-\(6\) 的 \(19\) 比特缺口只是这条普遍失配律在首个非平凡窗口上的精确实例。
\end{conjecture}

\noindent
其直接根据在于：命题 \ref{prop:fold-groupoid-z2x2-sector-by-orientation} 与定理 \ref{thm:op-algebra-fold-weak-hopf-z2x2-sector-closure} 已对一般 \(m\) 固定四扇区框架，而定理 \ref{thm:gut-fiber-parity-minimal-complete-bits} 在“全部纤维均非平凡”的情形下给出最小完备比特数恰为 \(|X_m|\)；因此随分辨率增长而线性放大的失配，至少在现有证书链上已呈现出稳定的全局形态。

\noindent
由此，MOM 前缀指数把四生成元投影词演算分裂为彼此隔离的乘法扇区；交叉异常族把轴可分解性的审计门槛精确提升为“全阶且不可降阶”的完备家族；对称幂商核则表明任意显影阈值仍可在多项式后端体积内实现；而当读出被限制为 window-\(6\) 的四个取向扇区时，完备异常所需的二十一比特信息被迫压缩为两比特粗标签。这条链把“可重写性”“可审计性”“可编译性”与“不可压缩性”锁进了同一组离散不变量。

\endinput


\endinput


\endinput
